\section{Discussion}
\label{sec:discussion}

The model developed in this paper is deliberately simple, but its implications are broader than a narrow automation problem. The central point is that when agentic AI is supplied by an external platform, the firm's organizational problem is not exhausted by current execution efficiency. Delegation to external AI changes the firm's internal authority structure because it also changes the stock of retained human recovery capability available when contingencies arise. In that sense, the relevant trade-off is not merely between human and machine performance on current tasks, but between present efficiency and future organizational resilience.

This perspective helps clarify why full automation need not be the natural limiting outcome of AI progress. In a purely static model, sufficiently capable AI would eventually displace human authority wherever it outperforms human handling. In the present framework, however, the value of human authority is not confined to contemporaneous task execution. Human authority also preserves internal recovery capability for intervention, override, and coordination under abnormal conditions, especially when externally supplied AI becomes unavailable or cannot be fully relied upon. The model therefore suggests that the persistence of human authority need not reflect technological backwardness. It may instead be an optimal organizational response to dependence on an external execution technology.

A related implication concerns the institutional role of platform pricing. In the model, pricing does not simply influence whether firms adopt AI. It also affects which decisions remain internal to the organization. This is important because it implies that the governance of external AI platforms has consequences inside the firm. A priced access regime changes the boundary between delegated and retained tasks, and therefore changes the scope of human discretion, oversight, and organizational learning. Because the pricing wedge is paid step by step, these distortions are stronger in settings where AI use is more genuinely agentic, that is, more reliant on long-horizon execution.

The model's notion of retained human recovery capability should be interpreted broadly as an internal experience stock relevant for abnormal states. It may represent crisis-response ability, override capacity, supervisory judgment, tacit operational knowledge, or the practical competence needed to diagnose and correct failures when the external platform cannot be fully relied upon. The reduced-form specification is intentionally parsimonious, but the underlying idea is general: internal authority sustains a stock of organizational capability that is valuable precisely when routine external execution becomes insufficient.

Several limitations are worth noting. First, the model takes platform pricing as exogenous. This is appropriate for isolating the firm's internal organization problem, but it leaves aside the upstream determination of access terms and the market structure of AI provision. Second, the model represents contingencies in reduced form rather than deriving them from a richer microfoundation of breakdown, override, or domain shift. Third, the firm is represented as a single decision-making unit, abstracting from strategic interaction across firms. These simplifications are deliberate, but they also indicate natural directions for future work.

Two extensions seem especially promising. One is to endogenize platform pricing or platform market power, thereby linking internal authority allocation to the industrial organization of AI provision. The other is to connect the present mechanism to make-or-buy decisions and firm boundaries. If external AI platforms perform routine execution while internal teams preserve recovery and exception-handling capability, then the trade-off identified here may also shape outsourcing, vertical scope, and the internal structure of the firm. These extensions lie beyond the present paper, but the baseline mechanism developed here provides a tractable starting point for them.
